\begin{table}[htbp]
\centering
\caption{\\ Stock market responses to other legal and political events}
\label{tab:other_events}
\footnotesize
\setlength{\tabcolsep}{4pt}
\medskip
\begin{adjustbox}{max width=\textwidth}
\begin{tabular}{lccccccc}
\toprule
 & & & & \multicolumn{2}{c}{Cumulative return (equal-wt.)} & & \\
\cmidrule(lr){5-6}
Event & Dates & Days & Market & Gold & No gold & Diff. & $t$ \\
\midrule
First gold-clause hearing (Irving Trust) & May 22--24, 1933 & 3 & +5.5\% & +8.3\% & +6.9\% & +1.4\% & 1.5 \\
\textit{In re Missouri Pacific} ruling & June 20--22, 1934 & 3 & -3.8\% & -5.1\% & -3.9\% & -1.2\% & -1.3 \\
\textit{Norman} affirmed (N.Y. Ct.\ App.) & July 3--6, 1934 & 3 & +2.2\% & +2.2\% & +1.5\% & +0.7\% & 0.7 \\
Cert.\ granted, \textit{Norman} & Oct.\ 8--10, 1934 & 3 & +0.7\% & +2.1\% & +1.4\% & +0.7\% & 0.7 \\
Cert.\ granted \& midterm election & Nov.\ 5--8, 1934 & 3 & +2.1\% & +7.5\% & +4.0\% & +3.5\% & 3.6 \\
\midrule
Joint Resolution (reference) & May 26--June 6, 1933 & 9 & +12.1\% & +30.9\% & +25.4\% & +5.5\% & 3.3 \\
\bottomrule
\end{tabular}
\end{adjustbox}
\medskip
\begin{minipage}{0.95\textwidth}\footnotesize \textit{Notes.} This table reports stock market responses to the intermediate legal and political events of the gold clause litigation, using the same portfolio construction as Table~\ref{tab:event_study} in the main text. Returns are cumulative raw returns, the compound product of daily returns over each event window. ``Diff.''\ is the difference between the cumulative returns of the gold-exposed and non-exposed equal-weighted portfolios. The $t$-statistic scales this difference by $\hat\sigma\sqrt{h}$, where $h$ is the number of trading days in the window and $\hat\sigma = 0.55$ percentage points is the standard deviation of the daily return differential between the two portfolios over calm trading days in 1934 (excluding the event windows in this table). The Irving Trust hearing (May~22--24, 1933) was the first court hearing on the enforceability of gold clauses. \textit{In re Missouri Pacific} (E.D.~Mo., June~20, 1934) was the first federal ruling upholding the constitutionality of the Joint Resolution; the New York Court of Appeals affirmed \textit{Norman v.\ Baltimore \& Ohio R.~Co.} on July~3, 1934. Certiorari was granted in \textit{Norman} on October~8, 1934 and in \textit{United States v.\ Bankers Trust Co.} on November~5, 1934; the New York Stock Exchange was closed on November~6, 1934 for the midterm election, so the November window mixes the certiorari grant with the election result. The Joint Resolution window from Table~\ref{tab:event_study} is repeated for reference.\end{minipage}
\end{table}
